%% Garcia's latex template for Activities Drescriptive Memorial Report
%% Version 0.2
%% (c) 2015 Vinicius Cardoso Garcia (vcg@cin.ufpe.br)
%%
%% This document is based on latex template from Prof. Daniel Cunha.
%%
%% Reference commands. Use the following commands to make references in your
%% text:
%%          \figref  -- for Figure reference
%%          \tabref  -- for Table reference
%%          \eqnref  -- for equation reference
%%          \chapref -- for chapter reference
%%          \secref  -- for section reference
%%          \appref  -- for appendix reference
%%          \axiref  -- for axiom reference
%%          \conjref -- for conjecture reference
%%          \defref  -- for definition reference
%%          \lemref  -- for lemma reference
%%          \theoref -- for theorem reference
%%          \corref  -- for corollary reference
%%          \propref -- for proprosition reference
%%          \pgref   -- for page reference
%%
%%          Example: See \chapref{chap:introduction}
%%%%%%%%%%%%%%%%%%%%%%%%%%%%%%%%%%%%%%%%%%%%%%%%%%%%%%%%%%%%%%%%%%%%%%%%%%%%%%%

\documentclass[a4paper,oneside,10pt]{article}

\usepackage{graphicx}
\usepackage{amsmath,amstext,amssymb,amsfonts}
\usepackage[none]{hyphenat}
\usepackage{fancyhdr}
\usepackage{cite}
\usepackage{indentfirst}
%\usepackage{path}
\usepackage[usenames, dvipsnames]{color}

\usepackage{pdfsync}
\usepackage{url}
\usepackage{setspace}
\usepackage[latin1,utf8]{inputenc}
\usepackage[T1]{fontenc}
\usepackage{colortbl}
\newcommand{\SetRowColor}[1]{\noalign{\gdef\RowColorName{#1}}\rowcolor{\RowColorName}}

\definecolor{MyRed}{rgb}{1,0.2,0.1}
\definecolor{light-gray}{gray}{0.95}
\definecolor{gray}{gray}{0.6}

\usepackage{booktabs}
\usepackage{ctable}
\usepackage{setspace}
\usepackage[multidot]{grffile}
\usepackage[final]{pdfpages}

\usepackage{titlesec}

\usepackage[toc,page]{appendix}
\newcommand{\appref}[1]{\@appendixname~\ref{#1}\xspace}
%\usepackage{xifthen}

\usepackage[bookmarks,colorlinks,pdfpagelabels,
pdftitle={Memorial Descritivo de Atividades}, pdfauthor={Vinicius Cardoso Garcia},
pdfsubject={Solicita\c{c}\~{a}o de progresss\~{a}o funcional docente de Adjunto N\'{\i}vel 3 para Adjunto N\'{\i}vel 4 apresentada \`{a} Comiss\~{a}o de Avalia\c{c}\~{a}o de Progress\~{a}o Horizontal do Centro de Inform\'{a}tica da Universidade Federal de Pernambuco.},
pdfcreator={Vinicius Cardoso Garcia}, pdfkeywords={Progressão, Adjunto, UFPE, CIn, Docente, Horizontal}]{hyperref}

\newcommand{\otoprule}{\midrule[\heavyrulewidth]}

% Defini\c{c}\~{a}o de margens
\setlength{\textwidth}{16cm} %
\setlength{\textheight}{23cm} %
\setlength{\oddsidemargin}{0cm} %
\setlength{\evensidemargin}{0cm} %
\setlength{\topmargin}{0cm}

\renewcommand{\abstractname}{Resumo}
\renewcommand{\contentsname}{\'{I}ndice Anal\'{\i}tico}
\renewcommand{\refname}{Refer\^{e}ncias}
\renewcommand{\appendixname}{Ap\^{e}ndice}
\renewcommand{\tablename}{Tabela}

%% Formata\c{c}\~{a}o do cabe\c{c}alho e rodap\'{e}
\lhead{\footnotesize Memorial Descritivo de Atividades} %
\rhead{\footnotesize \emph{Vinicius Cardoso Garcia}} %
\chead{} %
\cfoot{} %
\lfoot{\footnotesize\nouppercase\leftmark} %
\rfoot{\bfseries\thepage}
\renewcommand{\footrulewidth}{0.1pt}

% Comando para inserir n\'{u}mero de documento
\newcounter{document}%[section]
\setcounter{document}{0}
\renewcommand\theenumi{\arabic{section}.\arabic{enumi}}
\newcommand\Doc{{\addtocounter{document}{1}\mbox{\sffamily\bfseries [Doc. \arabic{document}]}}}

% Comando para repetir um n\'{u}mero de documento j\'{a} citado
% \mbox{\sffamily{\bfseries{[Doc. XX]}}}

%% Alternativa na edi\c{c}\~{a}o dos comandos.
%% Comando para inserir n\'{u}mero de documento
% \newcommand\thedocument{%
%    \ifthenelse{\arabic{subsection}=0}
%      {\thesection.\arabic{document}}
%      {\thesubsection.\arabic{document}}}
% \newcounter{document}[section]
% \setcounter{document}{0}
% \renewcommand\theenumi{\arabic{section}.\arabic{enumi}}
% \ifthenelse{\arabic{subsection} = 0}{\newcommand\Doc{{\stepcounter{document}\bfseries [Doc. \arabic{section}.\arabic{document}]}}}{\newcommand\Doc{{\stepcounter{document}\bfseries [Doc. \arabic{section}.\arabic{subsection}.\arabic{document}]}}}
% %\newcommand\Doc{{\addtocounter{document}{1}\mbox{\bfseries [Doc. \arabic{document}]}}}


% Ambiente para centralizar vertical
\newenvironment{vcenterpage}
     {\newpage\vspace*{\fill}}
     {\vspace*{\fill}\par\pagebreak}

\sloppy

\pagestyle{fancy}

\setcounter{secnumdepth}{4}

\begin{document}

\begin{titlepage}

\vspace{-5.0cm}

\begin{figure}[!htb]
 \centering{\includegraphics[scale=4.0]{logo2.eps}}
 \label{fig:UFPE_logo}
\end{figure}

\begin{center}
\vspace{1cm}
%{\huge \textsf{Solicita\c{c}\~{a}o de Progress\~{a}o Funcional Docente}} \\[1cm]
\rule{1.0\textwidth}{1pt} \\ [0.5cm]
{\Huge \textbf{\textsf{Memorial Descritivo de Atividades}}} \\
\rule{1.0\textwidth}{1pt} \\
\vspace{2cm}

\doublespacing
{\Large \textsf{Solicita\c{c}\~{a}o de progresss\~{a}o funcional docente de \textbf{Adjunto N\'{\i}vel 3 para Adjunto N\'{\i}vel 4} apresentada \`{a} Comiss\~{a}o de Avalia\c{c}\~{a}o de Progress\~{a}o Horizontal do Centro de Inform\'{a}tica da Universidade Federal de Pernambuco}}\\
\vspace{1.5cm}
{\LARGE \textsf{Solicitante: \textbf{Vinicius Cardoso Garcia}}}\\
\vspace{0.5cm}
{\Large \textsf{SIAPE: \textbf{1807586}}} \\
\vspace{0.5cm}
{\Large \textsf{Per\'{\i}odo: \textbf{20/08/2014 - 19/08/2016}}} \\

\vspace{2.0cm}

\normalsize \textsf{Setembro de 2016}

\end{center}
\thispagestyle{empty}
\end{titlepage}


\tableofcontents
%\include{Lista_Anexos} \cleartooddpage[\thispagestyle{empty}]

%%%%%%%%%%%%%%%%%%%%%%%%%%%%%%%%%%%%%%%%%%%%%%%%%%%%%%%%%%%%%%%%%%%%%%%%%%%%%%%
% APRESENTA\c{C}\~{A}O
%%%%%%%%%%%%%%%%%%%%%%%%%%%%%%%%%%%%%%%%%%%%%%%%%%%%%%%%%%%%%%%%%%%%%%%%%%%%%%%

\newpage
\section*{Apresenta\c{c}\~{a}o}
\vspace{0.3cm}

\begin{onehalfspace}

Memorial apresentado por \textbf{Vinicius Cardoso Garcia}, Professor Adjunto 3, lotado no Departamento de Informa\c{c}\~{a}o e Sistemas do Centro de Inform\'{a}tica (CIn) da Universidade Federal de Pernambuco (UFPE), para avalia\c{c}\~{a}o de desempenho acad\^{e}mico, para fins de acesso por Progress\~{a}o Horizontal \`{a} classe de Professor Adjunto N\'{\i}vel 4.

O presente memorial relata as atividades desempenhadas no per\'{\i}odo de \textbf{20 de Agosto de 2014 a 19 de Agosto de 2016} e foi elaborado com base nas diretrizes estabelecidas nas Resolu\c{c}\~{o}es 04/2008 e 02/2011 do Conselho Universit\'{a}rio (CONSUN) da UFPE e na Portaria 554, de 20/06/2013, do Minist\'{e}rio da Educa\c{c}\~{a}o (MEC). Os documentos comprobat\'{o}rios referenciados neste memorial est\~{a}o organizados em volumes anexos devidamente numerados. Por fim, os Anexos que se referem a artigos publicados em confer\^{e}ncias e peri\'{o}dicos cont\'{e}m apenas a primeira p\'{a}gina do trabalho e o email informando a aceita\c{c}\~{a}o do trabalhao para publica\c{c}\~{a}o (quando pertinente).

\end{onehalfspace}

\subsection*{Esclarecimentos}

Aqui nesta seção especial podem ser listadas algumas observa\c{c}\~{o}es acerca da documenta\c{c}\~{a}o comprobat\'{o}ria para dirimir poss\'{\i}veis d\'{u}vidas da Comiss\~{a}o Avaliadora. A numera\c{c}\~{a}o indicada corresponde \`{a} numera\c{c}\~{a}o do Anexo.

\begin{itemize}

\item [18.] Artigo dispon\'{\i}vel em \verb"dx.doi.org/10.14209/sbrt.2013.182".
\item [19.] Lista de artigos aceitos dispon\'{\i}vel em \verb"http://www.sbrt.org.br/sbrt2012/artigos.html". O artigo em quest\~{a}o possui identifica\c{c}\~{a}o 98879.
\item E assim por diante.

\end{itemize}


%%%%%%%%%%%%%%%%%%%%%%%%%%%%%%%%%%%%%%%%%%%%%%%%%%%%%%%%%%%%%%%%%%%%%%%%%%%%%%%
% Grupo 1 - Atividades de Ensino
%%%%%%%%%%%%%%%%%%%%%%%%%%%%%%%%%%%%%%%%%%%%%%%%%%%%%%%%%%%%%%%%%%%%%%%%%%%%%%%
\newpage
\section{Atividades de Ensino}

%A seguir, listo as atividades de ensino que realizei no per\'{\i}odo, separadas por subgrupo, conforme rege o documento tal.

%%%%%%%%%%%%%%%%%%%%%%%%%%%%%%%%%%%%%%%%%%%%%%%%%%%%%%%%%%%%%%%%%%%%%%%%%%%%%%%
% Subgrupo 1.1 - Orienta\c{c}\~{o}es e Co-Orienta\c{c}\~{o}es
%%%%%%%%%%%%%%%%%%%%%%%%%%%%%%%%%%%%%%%%%%%%%%%%%%%%%%%%%%%%%%%%%%%%%%%%%%%%%%%
\subsection{Orienta\c{c}\~{o}es e Co-Orienta\c{c}\~{o}es}
\vspace{0.3cm}

\subsubsection{Orienta\c{c}\~{a}o de Teses de Doutorado Conclu\'{i}das}
\vspace{0.3cm}

\begin{enumerate}
\renewcommand{\labelenumi}{{\large\bfseries\arabic{enumi}.}}

\item       \textbf{Aluno:} Jack Allan Bauer \mbox{\sffamily{\bfseries{[Doc. \ref{advisor:phd-concluidas}]}}} \\
            \textbf{T\'{\i}tulo da Tese:} Salvando o mundo em 24 horas.\\
            \textbf{Data da Defesa:} 6 de Novembro de 2001 \\
            \textbf{Institui\c{c}\~{a}o:} Universidade Federal de Pernambuco.

\end{enumerate}

%------------------------------------------------------------------------------

\subsubsection{Co-Orienta\c{c}\~{a}o de Teses de Doutorado Conclu\'{i}das}
\vspace{0.3cm}

\begin{enumerate}
\renewcommand{\labelenumi}{{\large\bfseries\arabic{enumi}.}}

\item       \textbf{Aluno:} Thomas A. Anderson \mbox{\sffamily{\bfseries{[Doc. \ref{co-advisor:phd-concluidas}]}}}\\
            \textbf{T\'{\i}tulo da Tese:} The Matrix \\
            \textbf{Data da Defesa:} 31 de Março de 1999 \\
            \textbf{Institui\c{c}\~{a}o:} Universidade Federal de Pernambuco.

\end{enumerate}

%------------------------------------------------------------------------------

\subsubsection{Orienta\c{c}\~{a}o de Teses de Doutorado em Andamento}
\vspace{0.3cm}

\begin{enumerate}
\renewcommand{\labelenumi}{{\large\bfseries\arabic{enumi}.}}

\item       \textbf{Aluno:} Morpheus Laurence Fishburne) \mbox{\sffamily{\bfseries{[Doc. \ref{advisor:phd-andamento}]}}} \\
            \textbf{T\'{\i}tulo da Tese:} Nabucodonozor in The Matrix \\
            \textbf{Data de In\'{\i}cio:} Mar\c{c}o de 2001 \\
            \textbf{Institui\c{c}\~{a}o:} CIn/UFPE

\end{enumerate}

%------------------------------------------------------------------------------

\subsubsection{Co-Orienta\c{c}\~{a}o de Teses de Doutorado em Andamento}
\vspace{0.3cm}

\begin{enumerate}
\renewcommand{\labelenumi}{{\large\bfseries\arabic{enumi}.}}

\item       \textbf{Aluno:} Anakin Skywalker \mbox{\sffamily{\bfseries{[Doc. \ref{co-advisor:phd-andamento}]}}}\\
            \textbf{T\'{\i}tulo da Tese:} A Ameaça Fantasma \\
            \textbf{Data de In\'{\i}cio:} Maio de 1999 \\
            \textbf{Institui\c{c}\~{a}o:} CIn/UFPE

\end{enumerate}

%------------------------------------------------------------------------------

\subsubsection{Orienta\c{c}\~{a}o de Disserta\c{c}\~{o}es de Mestrado Conclu\'{i}das}
\vspace{0.3cm}

\begin{enumerate}
\renewcommand{\labelenumi}{{\large\bfseries\arabic{enumi}.}}

\item       \textbf{Aluno:}  Obi-Wan Kenobi \mbox{\sffamily{\bfseries{[Doc. \ref{advisor:msc-concluidas}]}}} \\
            \textbf{T\'{\i}tulo da Disserta\c{c}\~{a}o:} A Vingança dos Sith\\
            \textbf{Tipo:} Acadêmico \\%| Profissional\\
            \textbf{Data da Defesa:} 19 de Maio de 2005\\
            \textbf{Institui\c{c}\~{a}o:} CIn/UFPE.

\item       \textbf{Aluno:} Darth Maul \mbox{\sffamily{\bfseries{[Doc. \ref{advisor:mprof-concluidas}]}}} \\
            \textbf{T\'{\i}tulo da Disserta\c{c}\~{a}o:} Lord of the Sith\\
            \textbf{Tipo:} Profissional\\
            \textbf{Data da Defesa:} 19 de Maio de 2005\\
            \textbf{Institui\c{c}\~{a}o:} CIn/UFPE.

\end{enumerate}

%------------------------------------------------------------------------------

\subsubsection{Co-Orienta\c{c}\~{a}o de Disserta\c{c}\~{o}es de Mestrado Conclu\'{i}das}
\vspace{0.3cm}

\begin{enumerate}
\renewcommand{\labelenumi}{{\large\bfseries\arabic{enumi}.}}

\item       \textbf{Aluno:} Anthony Edward Stark \mbox{\sffamily{\bfseries{[Doc. \ref{co-advisor:msc-concluidas}]}}} \\
            \textbf{T\'{\i}tulo da Disserta\c{c}\~{a}o:} Tales of Suspense\\
            \textbf{Tipo:} Acadêmico \\%| Profissional\\
            \textbf{Data da Defesa:} Março de 1963 \\
            \textbf{Institui\c{c}\~{a}o:} Marvel Comics

\item       \textbf{Aluno:} Henry Jonathan Pym \mbox{\sffamily{\bfseries{[Doc. \ref{co-advisor:mprof-concluidas}]}}} \\
            \textbf{T\'{\i}tulo da Disserta\c{c}\~{a}o:} Homem-Formiga\\
            \textbf{Tipo:} Profissional\\
            \textbf{Data da Defesa:} Março de 1963 \\
            \textbf{Institui\c{c}\~{a}o:} Marvel Comics

\end{enumerate}

%------------------------------------------------------------------------------

\subsubsection{Orienta\c{c}\~{a}o de Disserta\c{c}\~{o}es de Mestrado em Andamento}
\vspace{0.3cm}

\begin{enumerate}
\renewcommand{\labelenumi}{{\large\bfseries\arabic{enumi}.}}

\item       \textbf{Aluno:} Stanley Martin Lieber \mbox{\sffamily{\bfseries{[Doc. \ref{advisor:msc-andamento}]}}}\\
            \textbf{T\'{\i}tulo da Disserta\c{c}\~{a}o:} Stan Lee\\
            \textbf{Tipo:} Acadêmico \\%| Profissional\\
            \textbf{Data de In\'{\i}cio:} 22 de Dezembro de 1922\\
            \textbf{Institui\c{c}\~{a}o:} Marvel Comics

\item       \textbf{Aluno:} Robert Bruce Banner \mbox{\sffamily{\bfseries{[Doc. \ref{advisor:mprof-andamento}]}}}\\
            \textbf{T\'{\i}tulo da Disserta\c{c}\~{a}o:} Avengers vol. 1\\
            \textbf{Tipo:} Profissional\\
            \textbf{Data de In\'{\i}cio:} 22 de Dezembro de 1922\\
            \textbf{Institui\c{c}\~{a}o:} Marvel Comics
\end{enumerate}

%------------------------------------------------------------------------------
\subsubsection{Co-Orienta\c{c}\~{a}o de Disserta\c{c}\~{o}es de Mestrado em Andamento}
\vspace{0.3cm}

\begin{enumerate}
\renewcommand{\labelenumi}{{\large\bfseries\arabic{enumi}.}}

\item       \textbf{Aluno:} Erik Magnus Lehnsherr \mbox{\sffamily{\bfseries{[Doc. \ref{co-advisor:msc-andamento}]}}}\\
            \textbf{T\'{\i}tulo da Disserta\c{c}\~{a}o:} Magneto\\
            \textbf{Tipo:} Acadêmico \\%| Profissional\\
            \textbf{Data de In\'{\i}cio:} Setembro de 1963\\
            \textbf{Institui\c{c}\~{a}o:} Marvel Comics

\item       \textbf{Aluno:} Charles Francis Xavier \mbox{\sffamily{\bfseries{[Doc. \ref{co-advisor:mprof-andamento}]}}}\\
            \textbf{T\'{\i}tulo da Disserta\c{c}\~{a}o:} Magneto\\
            \textbf{Tipo:} Profissional\\
            \textbf{Data de In\'{\i}cio:} Setembro de 1963\\
            \textbf{Institui\c{c}\~{a}o:} Marvel Comics
\end{enumerate}

%------------------------------------------------------------------------------

\subsubsection{Orienta\c{c}\~{a}o de Trabalhos de Conclus\~{a}o de Curso}
\vspace{0.3cm}

\begin{enumerate}
\renewcommand{\labelenumi}{{\large\bfseries\arabic{enumi}.}}

\item       \textbf{Aluno:} Uhtred Ragnarson \mbox{\sffamily{\bfseries{[Doc. \ref{advisor:tcc}]}}}\\
            \textbf{Curso:} Crônicas Saxônicas\\
            \textbf{T\'{\i}tulo da Monografia:} O Último Reino\\
            \textbf{Data da Defesa:} Maio de 2004\\
            \textbf{Institui\c{c}\~{a}o:} Bebbanburg

\end{enumerate}

%------------------------------------------------------------------------------

\subsubsection{Orienta\c{c}\~{a}o de Monitoria}
\vspace{0.3cm}

\begin{enumerate}
\renewcommand{\labelenumi}{{\large\bfseries\arabic{enumi}.}}

\item   \textbf{Aluno:} Ragnar Lodbrok \mbox{\sffamily{\bfseries{[Doc. \ref{advisor:monitoria}]}}} \\
        \textbf{Disciplina:}  Engenharia Software e Sistemas (turma SI)\\
        \textbf{Curso:} Sistemas de Informação\\
        \textbf{Semestre:} 2015.2

\end{enumerate}

%------------------------------------------------------------------------------

\subsubsection{Orienta\c{c}\~{a}o de Trabalhos de Inicia\c{c}\~{a}o Cient\'{\i}fica}
\vspace{0.3cm}

\begin{enumerate}
\renewcommand{\labelenumi}{{\large\bfseries\arabic{enumi}.}}

\item   \textbf{Aluno:} Gorm den Gamle \mbox{\sffamily{\bfseries{[Doc. \ref{advisor:icc}]}}} \\
        \textbf{Projeto:} Pedras de Jelling \\
        \textbf{Tema:} Nórdico antigo \\
        \textbf{Per\'{\i}odo:} 936 - 938\\
        \textbf{Financiamento:} CAPES

\end{enumerate}

%------------------------------------------------------------------------------

\subsubsection{Orienta\c{c}\~{a}o de Trabalhos de Apoio Acad\^{e}mico}
\vspace{0.3cm}

\begin{enumerate}
\renewcommand{\labelenumi}{{\large\bfseries\arabic{enumi}.}}

\item   \textbf{Aluno:} Qui-Gon \mbox{\sffamily{\bfseries{[Doc. \ref{advisor:bia}]}}}\\
        \textbf{Projeto:} Uso da Internet das Coisas no acesso multimodal a informações em um Smart Campus.\\
        \textbf{Tema:} Comunica\c{c}\~{o}es sem Fio \\
        \textbf{Categoria:} Iniciação acadêmica (Graduando em Ciência da Computação) – Universidade Federal de Pernambuco.\\
        \textbf{Per\'{\i}odo:} 01 de Outubro de 2014 a 30 de Setembro de 2015 \\
        \textbf{Financiamento:} Programa Institucional de Bolsa de Incentivo Acadêmico – BIA, Fundação de Amparo à Ciência e Tecnologia do Estado de Pernambuco (FACEPE), processo número: BIA-0155-1.03/14.

\end{enumerate}

%%%%%%%%%%%%%%%%%%%%%%%%%%%%%%%%%%%%%%%%%%%%%%%%%%%%%%%%%%%%%%%%%%%%%%%%%%%%%%%
% Subgrupo 1.2 - Participa\c{c}\~{a}o em Comiss\~{o}es Examinadoras
%%%%%%%%%%%%%%%%%%%%%%%%%%%%%%%%%%%%%%%%%%%%%%%%%%%%%%%%%%%%%%%%%%%%%%%%%%%%%%%
\subsection{Participa\c{c}\~{a}o em Comiss\~{o}es Examinadoras}
\vspace{0.3cm}

\subsubsection{Bancas Examinadoras de Concurso}
\vspace{0.3cm}

\begin{enumerate}
\renewcommand{\labelenumi}{{\large\bfseries\arabic{enumi}.}}
\vspace{0.3cm}

\item       \textbf{Descrição:} Parecerista do(s) curso(s) de Sistemas de Informação da Avaliação de Cursos Superiores do Guia do Estudante (GE) 2015 \mbox{\sffamily{\bfseries{[Doc. \ref{app:2015-guia-estudante}]}}}\\
            \textbf{Instituição:} Guia do Estudante / Editora Abril \\
            \textbf{Data/Período:} 2015 %16 de Agosto de 2015

\end{enumerate}

%------------------------------------------------------------------------------

\subsubsection{Bancas Congressos de Inicia\c{c}\~{a}o Cient\'{i}fica ou de Extens\~{a}o}
\vspace{0.3cm}

\begin{enumerate}
\renewcommand{\labelenumi}{{\large\bfseries\arabic{enumi}.}}
\vspace{0.3cm}

\item       \textbf{Descrição:} Participação na Comissão Avaliadora dos trabalhos na área de Ciências Exatas apresentados na 19ª Jornada de Iniciação Científica da FACEPE \mbox{\sffamily{\bfseries{[Doc. \ref{app:2015-jic-facepe}]}}} \\
            \textbf{Instituição:} Fundação de Amparo a Ciência e Tecnologia do Estado de Pernambuco (FACEPE) \\
            \textbf{Data/Período:} 09 a 12 de Junho de 2015

\end{enumerate}

%------------------------------------------------------------------------------

\subsubsection{Bancas de Trabalho de Conclus\~{a}o de Curso}
\vspace{0.3cm}

\begin{enumerate}
\renewcommand{\labelenumi}{{\large\bfseries\arabic{enumi}.}}
\vspace{0.3cm}

\item       \textbf{Aluno:} José Welington de Almeida Filho \mbox{\sffamily{\bfseries{[Doc. \ref{app:bancas-tcc}]}}}\\
            \textbf{T\'{\i}tulo da Monografia:} Uma proposta de uso de Gamification para o ensino de software.\\
            \textbf{Curso:} Graduação em Ciência da Computação\\
            \textbf{Universidade:} Universidade Federal de Pernambuco (UFPE)\\
            \textbf{Data:} 28 de Julho de 2015\\

\end{enumerate}

%------------------------------------------------------------------------------

\subsubsection{Bancas de Disserta\c{c}\~{a}o de Mestrado}
\vspace{0.3cm}

\begin{enumerate}
\renewcommand{\labelenumi}{{\large\bfseries\arabic{enumi}.}}
\vspace{0.3cm}

\item       \textbf{Candidato:} Helaine Solange Lins Barreiros \mbox{\sffamily{\bfseries{[Doc. \ref{app:2015-msc-hslb}]}}} \\
            \textbf{T\'{\i}tulo da Disserta\c{c}\~{a}o:} Análise da Completude dos Relatos de Experimentos em Elasticidade na Computação em Nuvem: Um Mapeamento Sistemático\\
            \textbf{Tipo:} Acadêmico\\
            \textbf{Programa:} Programa de Pós-Graduação em Ciência da Computação\\
            \textbf{Universidade:} Universidade Federal de Pernambuco (UFPE)\\
            \textbf{Data:} 31 de Agosto de 2015

\end{enumerate}

%------------------------------------------------------------------------------

\subsubsection{Bancas de Qualifica\c{c}\~ao / Proposta de Tese de Doutorado}
\vspace{0.3cm}

\begin{enumerate}
\renewcommand{\labelenumi}{{\large\bfseries\arabic{enumi}.}}
\vspace{0.3cm}

\item       \textbf{Candidato:} Thiago Monteiro Prota \mbox{\sffamily{\bfseries{[Doc. \ref{app:2015-quali-phd-tmp}]}}} \\
            \textbf{T\'{\i}tulo da Qualifica\c{c}\~{a}o:} Análises Estrutural e Comportamental Orientadas a Conformidade para o Desenvolvimento de Aplicações Multimídia\\
            \textbf{Programa:} Programa de Pós-Graduação em Ciência da Computação\\
            \textbf{Universidade:} Universidade Federal de Pernambuco (UFPE)\\
            \textbf{Data:} 18 de Setembro de 2015

\end{enumerate}

%------------------------------------------------------------------------------

\subsubsection{Bancas de Tese de Doutorado}
%\vspace{0.3cm}

\begin{enumerate}
\renewcommand{\labelenumi}{{\large\bfseries\arabic{enumi}.}}
\vspace{0.3cm}

\item       \textbf{Candidato:} Julio Cesar Damasceno \mbox{\sffamily{\bfseries{[Doc. \ref{app:2015-phd-jcd}]}}} \\
            \textbf{T\'{\i}tulo da Disserta\c{c}\~{a}o:} UCloud: Uma Abordagem para Implantação de Nuvem Privada para a Administração Pública Federal\\
            \textbf{Programa:} Programa de Pós-Graduação em Ciência da Computação\\
            \textbf{Universidade:} Universidade Federal de Pernambuco (UFPE)\\
            \textbf{Data:} 03 de Setembro de 2015

\end{enumerate}

%------------------------------------------------------------------------------

\subsubsection{Participa\c{c}\~{a}o em Bancas de Sele\c{c}\~{a}o para ingresso e exames de Qualifica\c{c}\~ao de Programa de P\'{o}-Gradua\c{c}\~{a}o \textit{Stricto Sensu}}

\begin{enumerate}
\renewcommand{\labelenumi}{{\large\bfseries\arabic{enumi}.}}
\vspace{0.3cm}

\item       \textbf{Descrição:} Membro Externo da Comissão de Seleção para o Doutorado Sanduíche - PDSE da CAPES \mbox{\sffamily{\bfseries{[Doc. \ref{app:2015-pgcomp-pdse}]}}} \\
            \textbf{Programa:} Programa de Pós-Graduação em Ciência da Computação (PGCOMP) do aluno Alcemir Rodrigues Santos\\
            \textbf{Instituição:} Universidade Federal da Bahia (UFBA) \\
            \textbf{Per\'{\i}odo de Realiza\c{c}\~{a}o:} 15 de Junho de 2015

\end{enumerate}

%%%%%%%%%%%%%%%%%%%%%%%%%%%%%%%%%%%%%%%%%%%%%%%%%%%%%%%%%%%%%%%%%%%%%%%%%%%%%%%
% Subgrupo 1.3 - Atividades de Ensino na Graduação e na Pós-Graduação
%%%%%%%%%%%%%%%%%%%%%%%%%%%%%%%%%%%%%%%%%%%%%%%%%%%%%%%%%%%%%%%%%%%%%%%%%%%%%%%

\subsection{Atividades de Ensino na Graduação e na Pós-Graduação}
\vspace{0.3cm}

\subsubsection{Atividades de Ensino: Disciplinas na Gradua\c{c}\~{a}o}
\vspace{0.3cm}

Na Tabela \ref{Tab:Disc_Grad}, est\~{a}o listadas as disciplinas ministradas nos cursos de gradua\c{c}\~{a}o de Sistemas de Informa\c{c}\~{a}o no per\'{\i}odo de 20/08/2014 a 19/08/2016. As disciplinas est\~{a}o organizadas por semestre, com a identifica\c{c}\~{a}o do curso, da turma e da disciplina, esta \'{u}ltima com sua respectiva carga hor\'{a}ria. Adicionalmente, consta na referida Tabela a carga hor\'{a}ria m\'{e}dia semestral ao longo do per\'{\i}odo de avalia\c{c}\~{a}o.

\begin{table}[!htpb]
\centering \small
\caption{\texttt{Disciplinas de gradua\c{c}\~{a}o lecionadas no per\'{\i}odo considerado} }
\begin{tabular}{ccccc}
\toprule
\textbf{Semestre} & \textbf{Curso} & \textbf{Turma} & \textbf{Disciplina} & \textbf{C.H.} \\
\otoprule
  & Ci\^{e}ncia da Computa\c{c}\~{a}o & $I6$ & IF687 - Introdu\c{c}\~{a}o \`{a} Multim\'{\i}dia & 02\\
$2012.1$ & Ci\^{e}ncia da Computa\c{c}\~{a}o & $I6$ & IF690 - Hist\'{o}ria e Futuro da Computa\c{c}\~{a}o & 02\\
  & Engenharia da Computa\c{c}\~{a}o & $I6$ & IF690 - Hist\'{o}ria e Futuro da Computa\c{c}\~{a}o & 02\\
\cmidrule{1-5}
  & Ci\^{e}ncia da Computa\c{c}\~{a}o & $I6$ & IF687 - Introdu\c{c}\~{a}o \`{a} Multim\'{\i}dia & 02\\
$2012.2$ & Ci\^{e}ncia da Computa\c{c}\~{a}o & $I6$ & IF690 - Hist\'{o}ria e Futuro da Computa\c{c}\~{a}o & 02\\
  & Engenharia da Computa\c{c}\~{a}o & $E6$ & IF690 - Hist\'{o}ria e Futuro da Computa\c{c}\~{a}o & 02\\
  & Engenharia da Computa\c{c}\~{a}o & $E0$ & ES348 - Teoria da Informa\c{c}\~{a}o & 04\\
\cmidrule{1-5}
\multicolumn{4}{r}{\textbf{Carga Hor\'{a}ria M\'{e}dia Semestral} (2012)} & \textbf{08} \\
\otoprule
& Ci\^{e}ncia da Computa\c{c}\~{a}o & $I6$ & IF687 - Introdu\c{c}\~{a}o \`{a} Multim\'{\i}dia & 02\\
$2013.1$ & Ci\^{e}ncia da Computa\c{c}\~{a}o & $I6$ & IF690 - Hist\'{o}ria e Futuro da Computa\c{c}\~{a}o & 02\\
  & Engenharia da Computa\c{c}\~{a}o & $E6$ & IF690 - Hist\'{o}ria e Futuro da Computa\c{c}\~{a}o & 02\\
  & Engenharia da Computa\c{c}\~{a}o & $E0$ & IF856 - T\'{o}picos Avan\c{c}ados em Proc. Sinais & 04\\
\cmidrule{1-5}
& Ci\^{e}ncia da Computa\c{c}\~{a}o & $I6$ & IF687 - Introdu\c{c}\~{a}o \`{a} Multim\'{\i}dia & 02\\
$2013.2$ & Ci\^{e}ncia da Computa\c{c}\~{a}o & $I6$ & IF690 - Hist\'{o}ria e Futuro da Computa\c{c}\~{a}o & 02\\
  & Engenharia da Computa\c{c}\~{a}o & $I6$ & IF690 - Hist\'{o}ria e Futuro da Computa\c{c}\~{a}o & 02\\
\cmidrule{1-5}
\multicolumn{4}{r}{\textbf{Carga Hor\'{a}ria M\'{e}dia Semestral} (2013)} & \textbf{08} \\
\bottomrule
\end{tabular}
\label{Tab:Disc_Grad}
\end{table}

%%%%%%%%%%%%%%%%%%%%%%%%%%%%%%%%%%%%%%%%%%%%%%%%%%%%%%%%%%%%%%%%%%%%%%%%%%%%%%%
% Subgrupo 1.4 - Avaliação Didática Docente pelo Discente
%%%%%%%%%%%%%%%%%%%%%%%%%%%%%%%%%%%%%%%%%%%%%%%%%%%%%%%%%%%%%%%%%%%%%%%%%%%%%%%
\vspace{5.0cm}
\subsection{Avaliação Didática Docente pelo Discente}
\vspace{0.3cm}

%Nada a declarar neste subgrupo.

\begin{enumerate}
\renewcommand{\labelenumi}{{\large\bfseries\arabic{enumi}.}}

\item   \textbf{Período:} 2015.1 \textbf{[Doc. \ref{evaluation:2015-1}]}\\
        \textbf{Disciplina:} IF1006 - TÓPICOS AVANÇADOS EM SI 3, Turma: S0\\
        \textbf{Disciplina:} IF977 - ENGENHARIA DE SOFTWARE, Turma: S4\\
        \textbf{Nota:} 10.

\item   \textbf{Período:} 2014.2 \\
        \textbf{Disciplina:} IF1006 - TÓPICOS AVANÇADOS EM SI 3, Turma: S0\\
        \textbf{Disciplina:} IF977 - ENGENHARIA DE SOFTWARE, Turma: S4\\
        \textbf{Nota:} \textbf{Não houve avaliação didática docente pelo discente}.

\end{enumerate}

%%%%%%%%%%%%%%%%%%%%%%%%%%%%%%%%%%%%%%%%%%%%%%%%%%%%%%%%%%%%%%%%%%%%%%%%%%%%%%%
% Grupo 2: Atividades de Produ\c{c}\~{a}o Cient\'{\i}fica e T\'{e}cnica, Art\'{i}stica e Cultural
%%%%%%%%%%%%%%%%%%%%%%%%%%%%%%%%%%%%%%%%%%%%%%%%%%%%%%%%%%%%%%%%%%%%%%%%%%%%%%%
\newpage
\section{Atividades de Produ\c{c}\~{a}o Cient\'{\i}fica e T\'{e}cnica, Art\'{i}stica e Cultural}

%%%%%%%%%%%%%%%%%%%%%%%%%%%%%%%%%%%%%%%%%%%%%%%%%%%%%%%%%%%%%%%%%%%%%%%%%%%%%%%
% Subgrupo 2.1 - Produtividade de Pesquisa
%%%%%%%%%%%%%%%%%%%%%%%%%%%%%%%%%%%%%%%%%%%%%%%%%%%%%%%%%%%%%%%%%%%%%%%%%%%%%%%
\subsection{Produtividade de Pesquisa}
\vspace{0.3cm}

%------------------------------------------------------------------------------

\subsubsection{Bolsista de produtividade em pesquisa e em inova\c{c}\~{a}o tecnol\'{o}gica}
\vspace{0.3cm}

\begin{enumerate}
\renewcommand{\labelenumi}{{\large\bfseries\arabic{enumi}.}}

\item \textbf{T\'{\i}tulo do projeto:} The Death Star \textbf{[Doc. \ref{app:bolsista-prod}]}\\
      \textbf{N\'{u}mero do processo:} DV-1123-5.8/13\\
      \textbf{Financiador/Edital:} Conselho Nacional de Desenvolvimento Científico e Tecnológico, Edital MCTI/CNPq Nº 14/2013\\
      \textbf{Per\'{\i}odo (in\'{\i}cio-fim):} 2015 - 2018\\
      \textbf{Situação:} Desenvolvimento tecnológico\\
      \textbf{Natureza:} Pesquisa.

\end{enumerate}

%------------------------------------------------------------------------------

\subsubsection{Participa\c{c}\~{a}o em Eventos Cient\'{\i}ficos (com apresenta\c{c}\~{a}o de trabalho ou oferecimento de cursos, palestras ou debates)}
\vspace{0.3cm}

\begin{enumerate}
\renewcommand{\labelenumi}{{\large\bfseries\arabic{enumi}.}}

\item   \textbf{Evento:} Congresso Brasileiro de Software: Teoria e Prática \mbox{\sffamily{\bfseries{[Doc. \ref{app:2015-cbsoft}]}}} \\
        \textbf{Propósito:} Participante\\
        \textbf{Período:} 21 a 25 de Setembro de 2015\\
        \textbf{Local:} Belo Horizonte-MG, Brasil.

\item   \textbf{Evento:} XII Encontro Anual de Computação (EnAComp) \mbox{\sffamily{\bfseries{[Doc. \ref{app:2015-enacomp}]}}} \\
        \textbf{Propósito (i):} Ouvinte\\
        \textbf{Propósito (ii):} Apresentação da palestra ``Empreendedorismo e criação de negócios intensivos em software''\\
        \textbf{Propósito (ii):} Apresentação do minicurso ``Desenvolvimento de SaaS / Engenharia de Sofwtarte como Serviço'', com duração de 12 horas\\
        \textbf{Período:} 16 a 18 de Setembro de 2015\\
        \textbf{Local:} , Brasil.

\end{enumerate}

%------------------------------------------------------------------------------
% TODO
\subsubsection{Autoria de artigos completos publicados em anais de congresso, em jornais e revistas de circulação nacional e internacional na sua área}
\vspace{0.3cm}

\begin{enumerate}
\renewcommand{\labelenumi}{{\large\bfseries\arabic{enumi}.}}

\item GAMA, Kiev Santos; WANDERLEY, Rafael Roballo; MARANHÃO, Daniel Barbosa; GARCIA, Vinicius Cardoso. \textbf{TagHunt: Uma plataforma Combinando a Internet das Coisas com Scavenger Hunt Games}. In: 7º Simpósio Brasileiro de Computação Ubíqua e Pervasiva (SBCUP), Evento do XXXV Congresso da Sociedade Brasileira de Computação, Recife, 20 a 23 de Julho de 2015. \textbf{[Doc. \ref{conf:2015-sbcup}]}

\end{enumerate}

%------------------------------------------------------------------------------
% TODO
\subsubsection{Arbitragem de Artigos Técnico-Científicos Nacionais e Internacionais na sua área de atuação}
\vspace{0.3cm}

\begin{enumerate}
\renewcommand{\labelenumi}{{\large\bfseries\arabic{enumi}.}}

\item   \textbf{Peri\'{o}dico:} IEEE Transactions on Services Computing \textbf{[Doc. \ref{reviewer:2015-tsc}]}\\
        \textbf{Editora:} IEEE Computer Society\\
        \textbf{ISSN:} 1939-1374\\
        \textbf{URL:} \url{http://www.computer.org/web/tsc}

\end{enumerate}

%------------------------------------------------------------------------------
% TODO
\subsubsection{Coordena\c{c}\~{a}o e/ou Participa\c{c}\~{a}o em Projetos Aprovados por \'{O}rg\~{a}os de Fomento}
\vspace{0.3cm}

\begin{enumerate}
\renewcommand{\labelenumi}{{\large\bfseries\arabic{enumi}.}}

\item \textbf{T\'{\i}tulo do projeto:} BIGStore - Evolução da plataforma Ustore para Armazenamento, Manipulação e Experimentação de Grandes Volumes de Dados \textbf{[Doc. \ref{project:2015-facepe-pepe}]}\\
      \textbf{Fun\c{c}\~{a}o no projeto:} Integrante\\
      \textbf{N\'{u}mero do processo:} SIN-0199-1.03/15\\
      \textbf{Financiador/Edital:} Fundação de Amparo à Ciência e Tecnologia do Estado de Pernambuco, Edital FACEPE 23/2014 - PESQUISADOR NA EMPRESA DE PERNAMBUCO (PEPE)\\
      \textbf{Per\'{\i}odo (in\'{\i}cio-fim):} 2015 - 2018\\
      \textbf{Situação:} Desenvolvimento tecnológico\\
      \textbf{Natureza:} Pesquisa.

\end{enumerate}

%------------------------------------------------------------------------------

\subsubsection{Consultoria \`{a}s Institui\c{c}\~{o}es de Fomento \`{a} Pesquisa, Ensino e Extens\~{a}o}
\vspace{0.3cm}

\begin{enumerate}
\renewcommand{\labelenumi}{{\large\bfseries\arabic{enumi}.}}

\item   \textbf{Fun\c{c}\~{a}o:} Avaliador de Projetos de Pesquisa (modalidade Subven\c{c}\~{a}o Econ\^{o}mica - PAPPE Integra\c{c}\~{a}o) - Edital 10.2/2012 \textbf{[Doc. \ref{consulting:2015-facepe-pepe}]}\\
        \textbf{Institui\c{c}\~{a}o:} Fundação de Amparo à Ciência e Tecnologia do Estado de Pernambuco (FACEPE).

\end{enumerate}

%------------------------------------------------------------------------------

\subsubsection{Pr\^{e}mios Recebidos pela Produ\c{c}\~{a}o Cient\'{\i}fica e T\'{e}cnica}
\vspace{0.3cm}

Nada a declarar neste subgrupo.

\begin{enumerate}
\renewcommand{\labelenumi}{{\large\bfseries\arabic{enumi}.}}

\item Melhor artigo publicado no Simpósio Intergálito da República, 2014 \textbf{[Doc. \ref{award:2014}]}.

\end{enumerate}

%%%%%%%%%%%%%%%%%%%%%%%%%%%%%%%%%%%%%%%%%%%%%%%%%%%%%%%%%%%%%%%%%%%%%%%%%%%%%%%
% Subgrupo 2.2 - Produção Científica
%%%%%%%%%%%%%%%%%%%%%%%%%%%%%%%%%%%%%%%%%%%%%%%%%%%%%%%%%%%%%%%%%%%%%%%%%%%%%%%

\subsection{Produção Científica}
\vspace{0.3cm}

%------------------------------------------------------------------------------

\subsubsection{Trabalhos Publicados em Peri\'{o}dicos Especializados do Pa\'{\i}s ou do Exterior}
\vspace{0.3cm}

\begin{enumerate}
\renewcommand{\labelenumi}{{\large\bfseries\arabic{enumi}.}}

\item AFONSO, Ricardo Alexandre ; BRITO, Kellyton dos Santos ; DO NASCIMENTO, CLÓVIS HOLANDA ; COSTA, L.  C. ; ÁLVARO, ALEXANDRE ; Garcia, Vinicius Cardoso . ``(Br-SCMM) Brazilian Smart City Maturity Model: A Perspective from the Health Domain''. Studies in Health Technology and Informatics, v. 216 (MEDINFO 2015: eHealth-enabled Health), p. 983, 2015, doi: 10.3233/978-1-61499-564-7-983 \colorbox{yellow}{(Qualis XX)} \textbf{[Doc. \ref{journal:2015-medinfo}]}

\end{enumerate}

%%%%%%%%%%%%%%%%%%%%%%%%%%%%%%%%%%%%%%%%%%%%%%%%%%%%%%%%%%%%%%%%%%%%%%%%%%%%%%%
% Grupo 3: Atividades de Extens\~{a}o
%%%%%%%%%%%%%%%%%%%%%%%%%%%%%%%%%%%%%%%%%%%%%%%%%%%%%%%%%%%%%%%%%%%%%%%%%%%%%%%
\newpage
\section{Atividades de Extens\~{a}o}

%%%%%%%%%%%%%%%%%%%%%%%%%%%%%%%%%%%%%%%%%%%%%%%%%%%%%%%%%%%%%%%%%%%%%%%%%%%%%%%
% Subgrupo 3.1 - Coordenação e Orientação
%%%%%%%%%%%%%%%%%%%%%%%%%%%%%%%%%%%%%%%%%%%%%%%%%%%%%%%%%%%%%%%%%%%%%%%%%%%%%%%
\subsection{Subgrupo 3.1 - Coordenação e Orientação}
\vspace{0.3cm}

\begin{enumerate}
\renewcommand{\labelenumi}{{\large\bfseries\arabic{enumi}.}}

    \item Coordenação de Projeto de Extensão \textbf{[Doc. \ref{extensao:2015}]}

\end{enumerate}

%%%%%%%%%%%%%%%%%%%%%%%%%%%%%%%%%%%%%%%%%%%%%%%%%%%%%%%%%%%%%%%%%%%%%%%%%%%%%%%
% Subgrupo 3.2 - Coordenação de Eventos e Conferencista
%%%%%%%%%%%%%%%%%%%%%%%%%%%%%%%%%%%%%%%%%%%%%%%%%%%%%%%%%%%%%%%%%%%%%%%%%%%%%%%
\subsection{Coordenação de Eventos e Conferencista}
\vspace{0.3cm}

%------------------------------------------------------------------------------

\subsubsection{Comiss\~{a}o Organizadora de Eventos Internacional, Nacional, Regional ou Local}
\vspace{0.3cm}

\begin{enumerate}
\renewcommand{\labelenumi}{{\large\bfseries\arabic{enumi}.}}

    \item Program Committee Member of the 9th Brazilian Symposium on Software Components, Architectures and Reuse (SBCARS 2015) \textbf{[Doc. \ref{reviewer:2015-sbcars}]}

\end{enumerate}

%%%%%%%%%%%%%%%%%%%%%%%%%%%%%%%%%%%%%%%%%%%%%%%%%%%%%%%%%%%%%%%%%%%%%%%%%%%%%%%
% Grupo 4: Atividades de Forma\c{c}\~{a}o e Capacita\c{c}\~{a}o Acad\^{e}mica
%%%%%%%%%%%%%%%%%%%%%%%%%%%%%%%%%%%%%%%%%%%%%%%%%%%%%%%%%%%%%%%%%%%%%%%%%%%%%%%
\newpage
\section{Atividades de Forma\c{c}\~{a}o e Capacita\c{c}\~{a}o Acad\^{e}mica}
\vspace{0.3cm}

%------------------------------------------------------------------------------

\subsection{Atualiza\c{c}\~{a}o e Cursos de Capacita\c{c}\~{a}o ou Extens\~{a}o na \'{A}rea de Conhecimento ou Afins com no M\'{\i}nimo 40h}
\vspace{0.3cm}

\begin{enumerate}
\renewcommand{\labelenumi}{{\large\bfseries\arabic{enumi}.}}

\item   \textbf{Curso:}  Machine Learning (Coursera) \textbf{[Doc. \ref{courses:2015}]} \\
        \textbf{Modalidade:} Ensino \`{a} dist\^{a}ncia (EAD) \\
        \textbf{Carga Hor\'{a}ria:} 50 h  \\
        \textbf{Instrutor:} Prof. Andrew Ng \\
        \textbf{Institui\c{c}\~{a}o:} Stanford University

\end{enumerate}


%%%%%%%%%%%%%%%%%%%%%%%%%%%%%%%%%%%%%%%%%%%%%%%%%%%%%%%%%%%%%%%%%%%%%%%%%%%%%%%
% Grupo 5: Atividades Administrativas
%%%%%%%%%%%%%%%%%%%%%%%%%%%%%%%%%%%%%%%%%%%%%%%%%%%%%%%%%%%%%%%%%%%%%%%%%%%%%%%
\newpage
\section{Atividades Administrativas}
\vspace{0.3cm}

%------------------------------------------------------------------------------

\subsection{Membro de Comiss\~{a}o Tempor\'{a}ria}
\vspace{0.3cm}

\begin{enumerate}
\renewcommand{\labelenumi}{{\large\bfseries\arabic{enumi}.}}


\item   \textbf{Fun\c{c}\~{a}o:} Membro de comissão temporária \textbf{[Doc. \ref{committee:temp-2015}]}\\
        \textbf{Comiss\~{a}o:} Presidente\\
        \textbf{Per\'{\i}odo:} 01 a 15 de Outubro de 2015

\end{enumerate}

%------------------------------------------------------------------------------

\subsection{Coordenador de Curso Pós-Graduação \textbf{strictu sensu} }
\vspace{0.3cm}

\begin{enumerate}
\renewcommand{\labelenumi}{{\large\bfseries\arabic{enumi}.}}

\item   \textbf{Fun\c{c}\~{a}o:} Tutor da turma do Mestrado Profissional em Ciência da Computação com ênfase em Sistemas de Informação \textbf{[Doc. \ref{committee:postgrad-2015}]}\\
        \textbf{Per\'{\i}odo:} 01 de Dezembro de 2014 a 01 de Dezembro de 2017.

\end{enumerate}

%------------------------------------------------------------------------------

\subsection{Membro de Núcleo Docente Estruturante}
\vspace{0.3cm}

\begin{enumerate}
\renewcommand{\labelenumi}{{\large\bfseries\arabic{enumi}.}}

\item   \textbf{Fun\c{c}\~{a}o:} Membro do Núcleo Docente Estruturante do Curso de Graduação em Sistemas de Informação \textbf{[Doc. \ref{committee:nde-2015}]}\\
        \textbf{Per\'{\i}odo:} 21 de Julho de 2015 a 20 de Julho de 2016.

\item   \textbf{Fun\c{c}\~{a}o:} Membro do Núcleo Docente Estruturante do Curso de Graduação em Sistemas de Informação \textbf{[Doc. \ref{committee:nde-2014}]}\\
        \textbf{Per\'{\i}odo:} 21 de Julho de 2014 a 20 de Julho de 2015.

\end{enumerate}

%------------------------------------------------------------------------------

\subsection{Membro de Colegiados de Curso de Gradua\c{c}\~{a}o e P\'{o}s-Gradua\c{c}\~{a}o}
\vspace{0.3cm}

\begin{enumerate}
\renewcommand{\labelenumi}{{\large\bfseries\arabic{enumi}.}}

\item   \textbf{Fun\c{c}\~{a}o:} Membro do Colegiado da P\'{o}s-Gradua\c{c}\~{a}o \textbf{(em curso)}  \textbf{[Doc. \ref{committee:colegiado-postgrad-2015}]} \\
        \textbf{Per\'{\i}odo:} Desde Setembro de 2010.

\item   \textbf{Fun\c{c}\~{a}o:} Membro do Colegiado do Curso de Graduação em Ciência da Computação \textbf{[Doc. \ref{committee:colegiado-cc-2015}]} \\
        \textbf{Per\'{\i}odo:} 15 de Outubro de 2015 a 14 de Outubro de 2016.

\item   \textbf{Fun\c{c}\~{a}o:} Membro do Colegiado do Curso de Graduação em Sistemas de Informação \textbf{[Doc. \ref{committee:colegiado-si-2015}]} \\
        \textbf{Per\'{\i}odo:} 21 de Julho de 2015 a 20 de Julho de 2016.

\item   \textbf{Fun\c{c}\~{a}o:} Membro do Colegiado do Curso de Graduação em Sistemas de Informação \textbf{[Doc. \ref{committee:colegiado-si-2014}]} \\
        \textbf{Per\'{\i}odo:} 21 de Julho de 2014 a 20 de Julho de 2015.

\end{enumerate}

%%%%%%%%%%%%%%%%%%%%%%%%%%%%%%%%%%%%%%%%%%%%%%%%%%%%%%%%%%%%%%%%%%%%%%%%%%%%%%%
% LISTA DE ANEXOS
%%%%%%%%%%%%%%%%%%%%%%%%%%%%%%%%%%%%%%%%%%%%%%%%%%%%%%%%%%%%%%%%%%%%%%%%%%%%%%%

\newpage
\section{Lista de Anexos}

A Tabela \ref{Tab:ListaAnexos} cont\'{e}m a numera\c{c}\~{a}o e uma pequena descri\c{c}\~{a}o dos documentos comprobat\'{o}rios em anexo. Quando for o caso, h\'{a} uma indica\c{c}\~{a}o entre par\^{e}nteses, ao final da descri\c{c}\~{a}o, do setor respons\'{a}vel pela emiss\~{a}o do documento.


\begin{table}[h]
\small
\caption{\texttt{Rela\c{c}\~{a}o numerada dos Anexos comprobat\'{o}rios}.}
\begin{tabular}{cl}
\toprule
\large{\textbf{\texttt{Doc}}} & \multicolumn{1}{c}{\large{\textbf{\texttt{Descri\c{c}\~{a}o}}}} \\
\otoprule
  1 & Declara\c{c}\~{a}o de orienta\c{c}\~{a}o de Mestrado de Maxwell Silva (PPGES/UPE). \\
  %\cmidrule{1-2}
  2 & Declara\c{c}\~{a}o de orienta\c{c}\~{a}o de Mestrado de Bruna Melo (PPGES/UPE). \\
  %\cmidrule{1-2}
  3 & Declara\c{c}\~{a}o de orienta\c{c}\~{a}o de TCC de Diocleciano Neto e Lucas Paes (Gradua\c{c}\~{a}o EC CIn). \\
  %\cmidrule{1-2}
  4 & Termo de orienta\c{c}\~{a}o de Monitoria de Lucas Harada - 2013.2 (PROACAD/UFPE). \\
  %\cmidrule{1-2}
  5 & Declara\c{c}\~{a}o de orienta\c{c}\~{a}o de Monitoria de Lucas Harada - 2013.1 (Sec Grad CIn). \\
  %\cmidrule{1-2}
  6 & Declara\c{c}\~{a}o de orienta\c{c}\~{a}o de IC de Djeefther Albuquerque (PROPESQ/UFPE). \\
  %\cmidrule{1-2}
  7 & Declara\c{c}\~{a}o de participa\c{c}\~{a}o em banca de Doutorado de Isaac Benchimol (PPGEE/UFPE). \\
  %\cmidrule{1-2}
  8 & Declara\c{c}\~{a}o de participa\c{c}\~{a}o em banca de Doutorado de Walter Guimar\~{a}es (PPGEE/UFPE). \\
  \cmidrule{1-2}
  9 & Declara\c{c}\~{a}o de participa\c{c}\~{a}o em banca de Mestrado de Joyce Teixeira (PGCC/UFPE). \\
    & C\'{o}pia da Portaria PGCC 166/2012 - Composi\c{c}\~{a}o de banca de Mestrado (PGCC/UFPE). \\
  \cmidrule{1-2}
  10 & Declara\c{c}\~{a}o de participa\c{c}\~{a}o em banca de Mestrado de Antonio Assun\c{c}\~{a}o (PPGES/UPE). \\
     & C\'{o}pia da Portaria PPGES 11/2013 - Composi\c{c}\~{a}o de banca de Mestrado (PPGES/UPE). \\
  \cmidrule{1-2}
  11 & Declara\c{c}\~{a}o de participa\c{c}\~{a}o em banca de Mestrado de Maxwell Silva (PPGES/UPE). \\
     & C\'{o}pia da Portaria PPGES 15/2013 - Composi\c{c}\~{a}o de banca de Mestrado (PPGES/UPE). \\
  \cmidrule{1-2}
  12 & Declara\c{c}\~{a}o de participa\c{c}\~{a}o em banca de Mestrado de Bruna Melo (PPGES/UPE). \\
     & C\'{o}pia da Portaria PPGES 16/2013 - Composi\c{c}\~{a}o de banca de Mestrado (PPGES/UPE). \\
  \cmidrule{1-2}
  13 & Declara\c{c}\~{a}o de participa\c{c}\~{a}o em banca de TCC (Gradua\c{c}\~{a}o-EC CIn). \\
  %\cmidrule{1-2}
  14 & Declara\c{c}\~{a}o de particip. em banca de Sele\c{c}\~{a}o de Prof. Tempor\'{a}rio do CIn/UFPE (Dire\c{c}\~{a}o CIn). \\
  %\cmidrule{1-2}
  15 & Relat\'{o}rio de disciplinas lecionadas na gradua\c{c}\~{a}o no per\'{\i}odo de avalia\c{c}\~{a}o (SIGA-UFPE).\\
  %\cmidrule{1-2}
  16 & C\'{o}pia do certificado de participa\c{c}\~{a}o no XXX SBrT. \\
  %\cmidrule{1-2}
  17 & C\'{o}pia do certificado de participa\c{c}\~{a}o no XXXI SBrT. \\
  %\cmidrule{1-2}
  18 & C\'{o}pias do resumo publicado e do e-mail de aceita\c{c}\~{a}o de trabalho no XXXI SBrT. \\
  %\cmidrule{1-2}
  19 & C\'{o}pias do artigo publicado e do e-mail de aceita\c{c}\~{a}o de trabalho no XXX SBrT. \\
  %\cmidrule{1-2}
  20 a 23 & C\'{o}pias dos artigos publicados e dos e-mails de aceita\c{c}\~{a}o de trabalho no CBA 2012 e XXXI SBrT. \\
  %\cmidrule{1-2}
  24 a 31 & C\'{o}pias dos e-mails de convite para revis\~{a}o e finaliza\c{c}\~{a}o do processo (peri\'{o}dicos e eventos). \\
  \cmidrule{1-2}
  32 & C\'{o}pia do Resultado do Edital 10/2010 (pags 1,2 e 6) e Termo de Outorga do projeto (FACEPE). \\
     & C\'{o}pia de correspond\^{e}ncia comprobat\'{o}ria de presta\c{c}\~{a}o de contas do projeto (FACEPE). \\
  \cmidrule{1-2}
  33 & C\'{o}pias de e-mail de convite para avaliador e de cabe\c{c}alho do parecer emitido (FACEPE). \\
  %\cmidrule{1-2}
  34 & C\'{o}pia do certificado de premia\c{c}\~{a}o de melhor artigo de IC do XXXI SBrT. \\
  %\cmidrule{1-2}
  35 & C\'{o}pias do artigo publicado no peri\'{o}dico TEMA e do e-mail de aceita\c{c}\~{a}o. \\
  %\cmidrule{1-2}
  36 & C\'{o}pias do artigo publicado no peri\'{o}dico EURASIP J. Adv. Sig. Proc. e do e-mail de aceita\c{c}\~{a}o. \\
  %\cmidrule{1-2}
  37 & C\'{o}pias dos e-mails comprobat\'{o}rios de convite e resposta ao convite (MIC-CSC2012). \\
  %\cmidrule{1-2}
  38 & C\'{o}pias dos e-mails comprobat\'{o}rios de convite e resposta ao convite (XXXI SBrT). \\
  %\cmidrule{1-2}
  39 & C\'{o}pia da declara\c{c}\~{a}o enviada \`{a} Comiss\~{a}o Organizadora da SBPC 2013 (Dire\c{c}\~{a}o CIn). \\
  %\cmidrule{1-2}
  40 & Online Course Statement of Accomplishment: Machine Learning - Stanford University - Coursera. \\
  %\cmidrule{1-2}
  41 - 42 & Declara\c{c}\~{o}es de participa\c{c}\~{a}o na Comiss\~{a}o de Sele\c{c}\~{a}o do Mestrado CIn - 2012 e 2013. \\
  %\cmidrule{1-2}
  43 & Declara\c{c}\~{a}o comprobat\'{o}ria de cargo ocupado - Subcoordena\c{c}\~{a}o de Editais. \\
  %\cmidrule{1-2}
  44 & Declara\c{c}\~{a}o de participa\c{c}\~{a}o no Colegiado da P\'{o}s-Gradua\c{c}\~{a}o do CIn. \\
  %\cmidrule{1-2}
  45 & C\'{o}pia da Portaria de Designa\c{c}\~{a}o 005/2013 - Colegiado Gradua\c{c}\~{a}o EC CIn (Dire\c{c}\~{a}o CIn). \\
\bottomrule
\end{tabular}
\label{Tab:ListaAnexos}
\end{table}

% Appendix
\clearpage
%\addappheadtotoc
\appendix
%\appendixpage
\newpage
\section{Documentos comprobatórios}
Esta seção contém os documentos comprobatórios referentes às atividades listadas neste memorial.
\addcontentsline{toc}{section}{Documentos comprobatórios}
\renewcommand{\thesubsection}{\arabic{subsection}}
% \renewcommand{\subsection}{
% \titleformat{\subsection}
%   {\Huge\bfseries\center\vspace{.4\textwidth}\thispagestyle{fancy}} % format
%   {}                % label
%   {0pt}             % sep
%   {\huge}           % before-code
% }

%%%%%%%%%%%%%%%%%%%%%%%%%%%%%%%%%%%%%%%%%%%%%%%%%%%%%%%%%%%%%%%%%%%%%%%%%%%%%%%
% Grupo 1 - Atividades de Ensino
%%%%%%%%%%%%%%%%%%%%%%%%%%%%%%%%%%%%%%%%%%%%%%%%%%%%%%%%%%%%%%%%%%%%%%%%%%%%%%%

%%%%%%%%%%%%%%%%%%%%%%%%%%%%%%%%%%%%%%%%%%%%%%%%%%%%%%%%%%%%%%%%%%%%%%%%%%%%%%%
% Subgrupo 1.1 - Orienta\c{c}\~{o}es e Co-Orienta\c{c}\~{o}es
%%%%%%%%%%%%%%%%%%%%%%%%%%%%%%%%%%%%%%%%%%%%%%%%%%%%%%%%%%%%%%%%%%%%%%%%%%%%%%%

\newpage
\subsection{Orienta\c{c}\~{a}o de Teses de Doutorado Conclu\'{i}das}
\label{advisor:phd-concluidas}
Esta subseção apresenta o comprovante de orientações de Tese de Doutorado concluídas.
\includepdf[pages=-, scale=1,pagecommand=\thispagestyle{empty}]{\detokenize{GRUPO 1/Sub-Grupo 1.1/Comprovante Fake}}

\newpage
\subsection{Co-Orienta\c{c}\~{a}o de Teses de Doutorado Conclu\'{i}das}
\label{co-advisor:phd-concluidas}
Esta subseção apresenta o comprovante de co-orientações de Tese de Doutorado concluídas.
\includepdf[pages=-, scale=1,pagecommand=\thispagestyle{empty}]{\detokenize{GRUPO 1/Sub-Grupo 1.1/Comprovante Fake}}

\newpage
\subsection{Orienta\c{c}\~{a}o de Teses de Doutorado em Andamento}
\label{advisor:phd-andamento}
Esta subseção apresenta o comprovante de orientações de Tese de Doutorado em andamento.
\includepdf[pages=-, scale=1,pagecommand=\thispagestyle{empty}]{\detokenize{GRUPO 1/Sub-Grupo 1.1/Comprovante Fake}}

\newpage
\subsection{Co-Orienta\c{c}\~{a}o de Teses de Doutorado em Andamento}
\label{co-advisor:phd-andamento}
Esta subseção apresenta o comprovante de co-orientações de Tese de Doutorado em andamento.
\includepdf[pages=-, scale=1,pagecommand=\thispagestyle{empty}]{\detokenize{GRUPO 1/Sub-Grupo 1.1/Comprovante Fake}}

\newpage
\subsection{Orienta\c{c}\~{a}o de Disserta\c{c}\~{o}es de Mestrado Conclu\'{i}das}
\label{advisor:msc-concluidas}
Esta subseção apresenta o comprovante de orientações de Dissertação de Mestrado concluídas.
\includepdf[pages=-, scale=1,pagecommand=\thispagestyle{empty}]{\detokenize{GRUPO 1/Sub-Grupo 1.1/Comprovante Fake}}

\newpage
\subsection{Co-Orienta\c{c}\~{a}o de Disserta\c{c}\~{o}es de Mestrado Conclu\'{i}das}
\label{co-advisor:msc-concluidas}
Esta subseção apresenta o comprovante de co-orientações de Dissertação de Mestrado concluídas.
\includepdf[pages=-, scale=1,pagecommand=\thispagestyle{empty}]{\detokenize{GRUPO 1/Sub-Grupo 1.1/Comprovante Fake}}

\newpage
\subsection{Orienta\c{c}\~{a}o de Disserta\c{c}\~{o}es de Mestrado em Andamento}
\label{advisor:msc-andamento}
Esta subseção apresenta o comprovante de orientações de Dissertação de Mestrado em andamento.
\includepdf[pages=-, scale=1, pagecommand=\thispagestyle{empty}]{\detokenize{GRUPO 1/Sub-Grupo 1.1/Comprovante Fake}}

\newpage
\subsection{Co-Orienta\c{c}\~{a}o de Disserta\c{c}\~{o}es de Mestrado em Andamento}
\label{co-advisor:msc-andamento}
Esta subseção apresenta o comprovante de co-orientações de Dissertação de Mestrado em andamento.
\includepdf[pages=-, scale=1,pagecommand=\thispagestyle{empty}]{\detokenize{GRUPO 1/Sub-Grupo 1.1/Comprovante Fake}}

\newpage
\subsection{Orientação de Dissertação de Mestrado Profissional Conluídas}
\label{advisor:mprof-concluidas}
Esta subseção apresenta o comprovante de orientações de Dissertação de Mestrado Profissional concluídas.
\includepdf[pages=-, scale=1,pagecommand=\thispagestyle{empty}]{\detokenize{GRUPO 1/Sub-Grupo 1.1/Comprovante Fake}}

\newpage
\subsection{Co-Orientação de Dissertação de Mestrado Profissional Conluídas}
\label{co-advisor:mprof-concluidas}
Esta subseção apresenta o comprovante de co-orientações de Dissertação de Mestrado Profissional concluídas.
\includepdf[pages=-, scale=1,pagecommand=\thispagestyle{empty}]{\detokenize{GRUPO 1/Sub-Grupo 1.1/Comprovante Fake}}

\newpage
\subsection{Orientações de Mestrado Profissional em Andamento}
\label{advisor:mprof-andamento}
Esta subseção apresenta o comprovante de orientações de Dissertação de Mestrado Profissional em andamento.
\includepdf[pages=-, scale=1,pagecommand=\thispagestyle{empty}]{\detokenize{GRUPO 1/Sub-Grupo 1.1/Comprovante Fake}}

\newpage
\subsection{Co-Orientações de Mestrado Profissional em Andamento}
\label{co-advisor:mprof-andamento}
Esta subseção apresenta o comprovante de co-orientações de Dissertação de Mestrado Profissional em andamento.
\includepdf[pages=-, scale=1,pagecommand=\thispagestyle{empty}]{\detokenize{GRUPO 1/Sub-Grupo 1.1/Comprovante Fake}}

\newpage
\subsection{Orientações de Trabalhos de Conclusão de Curso Concluídas}
\label{advisor:tcc}
Esta subseção apresenta o comprovante de orientações de Trabalhos de Conclusão de Curso concluídas.
\includepdf[pages=-, scale=1,pagecommand=\thispagestyle{empty}]{\detokenize{GRUPO 1/Sub-Grupo 1.1/Comprovante Fake}}

\newpage
\subsection{Orientações de Monitorias}
\label{advisor:monitoria}
Esta subseção apresenta o comprovante de orientações de Monitoria no período de 2014-2 a 2016-1 impressa diretamente da página do Sistema de Gerenciamento de Monitores Online (GMon)\footnote{URL: \url{https://www.cin.ufpe.br/~gmon}, Último acesso em XX/XX/XXXX}.
\includepdf[pages=-, scale=1,pagecommand=\thispagestyle{empty}]{\detokenize{GRUPO 1/Sub-Grupo 1.1/Comprovante Fake}}

\newpage
\subsection{Orientações de Trabalhos de Apoio Acad\^{e}mico}
\label{advisor:icc}
Esta subseção apresenta o comprovante de orientações de Trabalhos de Iniciação Científica.
\includepdf[pages=-, scale=1,pagecommand=\thispagestyle{empty}]{\detokenize{GRUPO 1/Sub-Grupo 1.1/Comprovante Fake}}

\newpage
\subsection{Orientações de Trabalhos de Apoio Acad\^{e}mico}
\label{advisor:bia}
Esta subseção apresenta o comprovante de orientações de Trabalhos de Apoio Acad\^{e}mico.
\includepdf[pages=-, scale=1,pagecommand=\thispagestyle{empty}]{\detokenize{GRUPO 1/Sub-Grupo 1.1/Comprovante Fake}}


%%%%%%%%%%%%%%%%%%%%%%%%%%%%%%%%%%%%%%%%%%%%%%%%%%%%%%%%%%%%%%%%%%%%%%%%%%%%%%%
% Subgrupo 1.2 - Participa\c{c}\~{a}o em Comiss\~{o}es Examinadoras
%%%%%%%%%%%%%%%%%%%%%%%%%%%%%%%%%%%%%%%%%%%%%%%%%%%%%%%%%%%%%%%%%%%%%%%%%%%%%%%

\newpage
\subsection{Parecerista dos cursos de Sistemas de Informação da Avaliação de Cursos Superiores do Guia do Estudante (GE) 2015}
\label{app:2015-guia-estudante}
Esta subseção apresenta o comprovante da atuação como parecerista dos cursos de Sistemas de Informação da Avaliação de Cursos Superiores do Guia do Estudante (GE) 2015, promovido pela Editora Abril.
\includepdf[pages=-, scale=1,pagecommand=\thispagestyle{empty}]{\detokenize{GRUPO 1/Sub-Grupo 1.2/1/Comprovante Fake}}

%---

\newpage
\subsection{Comissão Avaliadora dos trabalhos na área de Ciências Exatas apresentados na 19ª Jornada de Iniciação Científica da FACEPE}
\label{app:2015-jic-facepe}
Esta subseção apresenta o comprovante da atuação como Membro da Comissão Avaliadora dos trabalhos na área de Ciências Exatas apresentados na 19ª Jornada de Iniciação Científica da Fundação de Amparo a Ciência e Tecnologia do Estado de Pernambuco (FACEPE).
\includepdf[pages=-, scale=1,pagecommand=\thispagestyle{empty}]{\detokenize{GRUPO 1/Sub-Grupo 1.2/2/Comprovante Fake}}

%---

\newpage
\subsection{Participação em Bancas de Trabalho de Conclus\~{a}o de Curso}
\label{app:bancas-tcc}
Esta subseção apresenta o comprovante da atuação membro nas Bancas Examinadoras de Trabalho de Conclusão de Curso no Centro de Informática/UFPE.
\includepdf[pages=-, scale=1,pagecommand=\thispagestyle{empty}]{\detokenize{GRUPO 1/Sub-Grupo 1.2/4/Comprovante Fake}}

\newpage
\subsection{Participação em Banca de Disserta\c{c}\~{a}o de Mestrado de Helaine Solange Lins Barreiros}
\label{app:2015-msc-hslb}
Esta subseção apresenta o comprovante da atuação membro na Banca Examinadora de Defesa de Disserta\c{c}\~{a}o de Mestrado de Helaine Solange Lins Barreiros no Centro de Informática/UFPE.
\includepdf[pages=-, scale=1,pagecommand=\thispagestyle{empty}]{\detokenize{GRUPO 1/Sub-Grupo 1.2/5/Comprovante Fake}}

%---

\newpage
\subsection{Participação em Banca de Qualifica\c{c}\~ao / Proposta de Tese de Doutorado de Thiago Monteiro Prota.}
\label{app:2015-quali-phd-tmp}
Esta subseção apresenta o comprovante da atuação como membro na Banca Examinadora de Defesa de Proposta de Tese de Doutorado de Thiago Monteiro Prota no Centro de Informática/UFPE.
\includepdf[pages=-, scale=1,pagecommand=\thispagestyle{empty}]{\detokenize{GRUPO 1/Sub-Grupo 1.2/6/Comprovante Fake}}

%---

\newpage
\subsection{Participação em Banca de Tese de Doutorado de Julio Cesar Damasceno.}
\label{app:2015-phd-jcd}
Esta subseção apresenta o comprovante da atuação como membro na Banca Examinadora de Defesa de Tese de Doutorado de Julio Cesar Damasceno no Centro de Informática/UFPE.
\includepdf[pages=-, scale=1,pagecommand=\thispagestyle{empty}]{\detokenize{GRUPO 1/Sub-Grupo 1.2/7/Comprovante Fake}}

%---

\newpage
\subsection{Participa\c{c}\~{a}o como Membro Externo da Comissão de Seleção para o Doutorado Sanduíche - PDSE da CAPES}
\label{app:2015-pgcomp-pdse}
Esta subseção apresenta o comprovante da atuação como membro Externo da Comissão de Seleção para o Doutorado Sanduíche - PDSE da CAPES.
\includepdf[pages=-, scale=1,pagecommand=\thispagestyle{empty}]{\detokenize{GRUPO 1/Sub-Grupo 1.2/8/Comprovante Fake}}

%%%%%%%%%%%%%%%%%%%%%%%%%%%%%%%%%%%%%%%%%%%%%%%%%%%%%%%%%%%%%%%%%%%%%%%%%%%%%%%
% Subgrupo 1.3 - Atividades de Ensino na Graduação e na Pós-Graduação
%%%%%%%%%%%%%%%%%%%%%%%%%%%%%%%%%%%%%%%%%%%%%%%%%%%%%%%%%%%%%%%%%%%%%%%%%%%%%%%



%%%%%%%%%%%%%%%%%%%%%%%%%%%%%%%%%%%%%%%%%%%%%%%%%%%%%%%%%%%%%%%%%%%%%%%%%%%%%%%
% Subgrupo 1.4 - Avaliação Didática Docente pelo Discente
%%%%%%%%%%%%%%%%%%%%%%%%%%%%%%%%%%%%%%%%%%%%%%%%%%%%%%%%%%%%%%%%%%%%%%%%%%%%%%%

\newpage
\subsection{Avaliação Didática Docente pelo Discente}
\label{evaluation:2015-1}
Esta subseção apresenta o comprovante da Avaliação Didática Docente pelo Discente para o período do primeiro semestre de 2015.
\includepdf[pages=-, scale=1,pagecommand=\thispagestyle{empty}]{\detokenize{GRUPO 1/Sub-Grupo 1.4/Comprovante Fake}}

%%%%%%%%%%%%%%%%%%%%%%%%%%%%%%%%%%%%%%%%%%%%%%%%%%%%%%%%%%%%%%%%%%%%%%%%%%%%%%%
% Subgrupo 2.1 - Produtividade de Pesquisa
%%%%%%%%%%%%%%%%%%%%%%%%%%%%%%%%%%%%%%%%%%%%%%%%%%%%%%%%%%%%%%%%%%%%%%%%%%%%%%%

\newpage
\subsection{Bolsista de produtividade em pesquisa e em inova\c{c}\~{a}o tecnol\'{o}gica}
\label{app:bolsista-prod}
Esta subseção apresenta o comprovante de Bolsista de produtividade em pesquisa e em inova\c{c}\~{a}o tecnol\'{o}gica.
\includepdf[pages=-, scale=1,pagecommand=\thispagestyle{empty}]{\detokenize{GRUPO 2/Sub-Grupo 2.1/A/Comprovante Fake}}

\newpage
\subsection{Participa\c{c}\~{a}o em Eventos Cient\'{\i}ficos (com apresenta\c{c}\~{a}o de trabalho ou oferecimento de cursos, palestras ou debates}
\label{app:2015-cbsoft}
Esta subseção apresenta o comprovante da participação no Congresso Brasileiro de Software: Teoria e Prática (CBSoft'2015) com seus respectivos propósitos.
\includepdf[pages=-, scale=1,pagecommand=\thispagestyle{empty}]{\detokenize{GRUPO 2/Sub-Grupo 2.1/B/Comprovante Fake}}

\newpage
\subsection{Participa\c{c}\~{a}o em Eventos Cient\'{\i}ficos (com apresenta\c{c}\~{a}o de trabalho ou oferecimento de cursos, palestras ou debates}
\label{app:2015-enacomp}
Esta subseção apresenta o comprovante da participação no XII Encontro Anual de Computação (EnAComp) com seus respectivos propósitos.
\includepdf[pages=-, scale=1,pagecommand=\thispagestyle{empty}]{\detokenize{GRUPO 2/Sub-Grupo 2.1/B/Comprovante Fake}}
\includepdf[pages=-, scale=1,pagecommand=\thispagestyle{empty}]{\detokenize{GRUPO 2/Sub-Grupo 2.1/B/Comprovante Fake}}
\includepdf[pages=-, scale=1,pagecommand=\thispagestyle{empty}]{\detokenize{GRUPO 2/Sub-Grupo 2.1/B/Comprovante Fake}}

%---

\newpage
\subsection{Autoria de artigos completos publicados em anais de congresso, em jornais e revistas de circulação nacional e internacional na sua área}
\label{conf:2015-sbcup}
Esta subseção apresenta o comprovante da autoria de artigo completo publicado em anais do 7º Simpósio Brasileiro de Computação Ubíqua e Pervasiva (SBCUP)\footnote{\url{http://csbc2015.cin.ufpe.br/eventos_descricao/13}}, Evento do XXXV Congresso da Sociedade Brasileira de Computação.
\includepdf[pages=-, scale=1,pagecommand=\thispagestyle{empty}]{\detokenize{GRUPO 2/Sub-Grupo 2.1/C/Comprovante Fake}}

%---

\newpage
\subsection{Arbitragem de Artigos Técnico-Científicos Nacionais e Internacionais na sua área de atuação}
\label{reviewer:2015-tsc}
Esta subseção apresenta o comprovante de arbitragem de artigos do periódico IEEE Transactions on Services Computing, editora IEEE Computer Society, ISSN: 1939-1374.
\includepdf[pages=-, scale=1,pagecommand=\thispagestyle{empty}]{\detokenize{GRUPO 2/Sub-Grupo 2.1/D/Comprovante Fake}}

%---

\newpage
\subsection{Coordena\c{c}\~{a}o e/ou Participa\c{c}\~{a}o em Projetos Aprovados por \'{O}rg\~{a}os de Fomento}
\label{project:2015-facepe-pepe}
Esta subseção apresenta o comprovante de coordena\c{c}\~{a}o e/ou participa\c{c}\~{a}o no Projeto \textit{``BIGStore - Evolução da plataforma Ustore para Armazenamento, Manipulação e Experimentação de Grandes Volumes de Dados''} aprovado pela Fundação de Amparo à Ciência e Tecnologia do Estado de Pernambuco (FACEPE), Edital FACEPE 23/2014 - PESQUISADOR NA EMPRESA DE PERNAMBUCO (PEPE).
\includepdf[pages=-, scale=1,pagecommand=\thispagestyle{empty}]{\detokenize{GRUPO 2/Sub-Grupo 2.1/E/Comprovante Fake}}

%---

\newpage
\subsection{Consultoria \`{a}s Institui\c{c}\~{o}es de Fomento \`{a} Pesquisa, Ensino e Extens\~{a}o}
\label{consulting:2015-facepe-pepe}
Esta subseção apresenta o comprovante de consultoria \`{a} Fundação de Amparo à Ciência e Tecnologia do Estado de Pernambuco (FACEPE) no Seminário de Avaliação Final dos Projetos Aprovados no Edital 10.2/2012 - PAPPE INTEGRAÇÃO 04a RODADA.
\includepdf[pages=-, scale=1,pagecommand=\thispagestyle{empty}]{\detokenize{GRUPO 2/Sub-Grupo 2.1/F/Comprovante Fake}}

%---

\newpage
\subsection{Pr\^{e}mios Recebidos pela Produ\c{c}\~{a}o Cient\'{\i}fica e T\'{e}cnica}
\label{award:2014}
Esta subseção apresenta o comprovante de consultoria \`{a} Fundação de Amparo à Ciência e Tecnologia do Estado de Pernambuco (FACEPE) no Seminário de Avaliação Final dos Projetos Aprovados no Edital 10.2/2012 - PAPPE INTEGRAÇÃO 04a RODADA.
\includepdf[pages=-, scale=1,pagecommand=\thispagestyle{empty}]{\detokenize{GRUPO 2/Sub-Grupo 2.1/G/Comprovante Fake}}

%%%%%%%%%%%%%%%%%%%%%%%%%%%%%%%%%%%%%%%%%%%%%%%%%%%%%%%%%%%%%%%%%%%%%%%%%%%%%%%
% Subgrupo 2.2 - Produção Científica
%%%%%%%%%%%%%%%%%%%%%%%%%%%%%%%%%%%%%%%%%%%%%%%%%%%%%%%%%%%%%%%%%%%%%%%%%%%%%%%

\newpage
\subsection{Trabalhos Publicados em Peri\'{o}dicos Especializados do Pa\'{\i}s ou do Exterior}
\label{journal:2015-medinfo}
Esta subseção apresenta o comprovante de publicação do artigo ```(Br-SCMM) Brazilian Smart City Maturity Model: A Perspective from the Health Domain'' no Studies in Health Technology and Informatics, v. 216 (MEDINFO 2015: eHealth-enabled Health), p. 983, 2015, doi: 10.3233/978-1-61499-564-7-983.
\includepdf[pages=-, scale=1,pagecommand=\thispagestyle{empty}]{\detokenize{GRUPO 2/Sub-Grupo 2.2/Comprovante Fake}}

%%%%%%%%%%%%%%%%%%%%%%%%%%%%%%%%%%%%%%%%%%%%%%%%%%%%%%%%%%%%%%%%%%%%%%%%%%%%%%%
% Grupo 3: Atividades de Extens\~{a}o
%%%%%%%%%%%%%%%%%%%%%%%%%%%%%%%%%%%%%%%%%%%%%%%%%%%%%%%%%%%%%%%%%%%%%%%%%%%%%%%

%%%%%%%%%%%%%%%%%%%%%%%%%%%%%%%%%%%%%%%%%%%%%%%%%%%%%%%%%%%%%%%%%%%%%%%%%%%%%%%
% Subgrupo 3.1 - Coordenação e Orientação
%%%%%%%%%%%%%%%%%%%%%%%%%%%%%%%%%%%%%%%%%%%%%%%%%%%%%%%%%%%%%%%%%%%%%%%%%%%%%%%

\newpage
\subsection{Coordenação e Orientação de Atividades de Extens\~{a}o}
\label{extensao:2015}
Esta subseção apresenta o comprovante de Coordenação e Orientação de Atividades de Extens\~{a}o.
\includepdf[pages=-, scale=1,pagecommand=\thispagestyle{empty}]{\detokenize{GRUPO 3/Sub-Grupo 3.1/Comprovante Fake}}

%%%%%%%%%%%%%%%%%%%%%%%%%%%%%%%%%%%%%%%%%%%%%%%%%%%%%%%%%%%%%%%%%%%%%%%%%%%%%%%
% Subgrupo 3.2 - Coordenação de Eventos e Conferencista
%%%%%%%%%%%%%%%%%%%%%%%%%%%%%%%%%%%%%%%%%%%%%%%%%%%%%%%%%%%%%%%%%%%%%%%%%%%%%%%

\newpage
\subsection{Comiss\~{a}o Organizadora de Eventos Internacional, Nacional, Regional ou Local}
\label{reviewer:2015-sbcars}
Esta subseção apresenta o comprovante de Comiss\~{a}o Organizadora do 9th Brazilian Symposium on Software Components, Architectures and Reuse (SBCARS 2015).
\includepdf[pages=-, scale=1,pagecommand=\thispagestyle{empty}]{\detokenize{GRUPO 3/Sub-Grupo 3.2/Comprovante Fake}}

%%%%%%%%%%%%%%%%%%%%%%%%%%%%%%%%%%%%%%%%%%%%%%%%%%%%%%%%%%%%%%%%%%%%%%%%%%%%%%%
% Grupo 4: Atividades de Forma\c{c}\~{a}o e Capacita\c{c}\~{a}o Acad\^{e}mica
%%%%%%%%%%%%%%%%%%%%%%%%%%%%%%%%%%%%%%%%%%%%%%%%%%%%%%%%%%%%%%%%%%%%%%%%%%%%%%%

\newpage
\subsection{Atualiza\c{c}\~{a}o e Cursos de Capacita\c{c}\~{a}o ou Extens\~{a}o na \'{A}rea de Conhecimento ou Afins com no M\'{\i}nimo 40h}
\label{courses:2015}
Esta subseção apresenta o comprovante de Atualiza\c{c}\~{a}o e Cursos de Capacita\c{c}\~{a}o ou Extens\~{a}o no curso Machine Learning (Coursera).
\includepdf[pages=-, scale=1,pagecommand=\thispagestyle{empty}]{\detokenize{GRUPO 4/Comprovante Fake}}

%%%%%%%%%%%%%%%%%%%%%%%%%%%%%%%%%%%%%%%%%%%%%%%%%%%%%%%%%%%%%%%%%%%%%%%%%%%%%%%
% Grupo 5: Atividades Administrativas
%%%%%%%%%%%%%%%%%%%%%%%%%%%%%%%%%%%%%%%%%%%%%%%%%%%%%%%%%%%%%%%%%%%%%%%%%%%%%%%

\newpage
\subsection{Atividades Administrativas, Membro de Comiss\~{a}o Tempor\'{a}ria}
\label{committee:temp-2015}
Esta subseção apresenta o comprovante de Atividades Administrativas, Membro de Comiss\~{a}o Tempor\'{a}ria.
\includepdf[pages=-, scale=1,pagecommand=\thispagestyle{empty}]{\detokenize{GRUPO 5/Comprovante Fake}}

%------------------------------------------------------------------------------

\newpage
\subsection{Atividades Administrativas, Coordenador de Curso Pós-Graduação \textbf{strictu sensu}}
\label{committee:postgrad-2015}
Esta subseção apresenta o comprovante de Atividades Administrativas, Coordenador de Curso Pós-Graduação \textbf{strictu sensu}.
\includepdf[pages=-, scale=1,pagecommand=\thispagestyle{empty}]{\detokenize{GRUPO 5/Comprovante Fake}}

%------------------------------------------------------------------------------

\newpage
\subsection{Atividades Administrativas, Membro de Núcleo Docente Estruturante}
\label{committee:nde-2015}
Esta subseção apresenta o comprovante de Atividades Administrativas, Membro de Núcleo Docente Estruturante.
\includepdf[pages=-, scale=1,pagecommand=\thispagestyle{empty}]{\detokenize{GRUPO 5/Comprovante Fake}}

\newpage
\subsection{Atividades Administrativas, Membro de Núcleo Docente Estruturante}
\label{committee:nde-2014}
Esta subseção apresenta o comprovante de Atividades Administrativas, Membro de Núcleo Docente Estruturante.
\includepdf[pages=-, scale=1,pagecommand=\thispagestyle{empty}]{\detokenize{GRUPO 5/Comprovante Fake}}

%------------------------------------------------------------------------------

\newpage
\subsection{Atividades Administrativas, Membro de Colegiados de Curso de Gradua\c{c}\~{a}o e P\'{o}s-Gradua\c{c}\~{a}o}
\label{committee:colegiado-postgrad-2015}
Esta subseção apresenta o comprovante de Atividades Administrativas, Membro do Colegiado da P\'{o}s-Gradua\c{c}\~{a}o (em curso).
\includepdf[pages=-, scale=1,pagecommand=\thispagestyle{empty}]{\detokenize{GRUPO 5/Comprovante Fake}}

\newpage
\subsection{Atividades Administrativas, Membro de Colegiados de Curso de Gradua\c{c}\~{a}o e P\'{o}s-Gradua\c{c}\~{a}o}
\label{committee:colegiado-cc-2015}
Esta subseção apresenta o comprovante de Atividades Administrativas, Membro do Colegiado do Curso de Graduação em Ciência da Computação.
\includepdf[pages=-, scale=1,pagecommand=\thispagestyle{empty}]{\detokenize{GRUPO 5/Comprovante Fake}}

\newpage
\subsection{Atividades Administrativas, Membro de Colegiados de Curso de Gradua\c{c}\~{a}o e P\'{o}s-Gradua\c{c}\~{a}o}
\label{committee:colegiado-si-2015}
Esta subseção apresenta o comprovante de Atividades Administrativas, Membro do Colegiado do Curso de Graduação em Sistemas de Informação em 2015.
\includepdf[pages=-, scale=1,pagecommand=\thispagestyle{empty}]{\detokenize{GRUPO 5/Comprovante Fake}}

\newpage
\subsection{Atividades Administrativas, Membro de Colegiados de Curso de Gradua\c{c}\~{a}o e P\'{o}s-Gradua\c{c}\~{a}o}
\label{committee:colegiado-si-2014}
Esta subseção apresenta o comprovante de Atividades Administrativas, Membro do Colegiado do Curso de Graduação em Sistemas de Informação em 2014.
\includepdf[pages=-, scale=1,pagecommand=\thispagestyle{empty}]{\detokenize{GRUPO 5/Comprovante Fake}}

%%%%%%%%%%%%%%%%%%%%%%%%%%%%%%%%%%%%%%%%%%%%%%%%%%%%%%%%%%%%%%%%%%%%%%%%%%%%%%%
% \newpage
% \subsection{Portaria de Progressão}
% \label{app:2014-portaria-progressao}
% \includepdf[pages=-, scale=1,pagecommand=\thispagestyle{empty}]{\detokenize{GRUPO 1/20141205_Portaria-de-Progressao-Funcional_5929-2014.pdf}}


\end{document}

%%% EOF
